\label{sec:conclusion}
\hspace*{1cm} Энэхүү бакалаврын судалгааны ажлын хүрээнд шинээр хөгжүүлэгдэж буй МУИС-ийн Дипломын ажлын удирдлагын системийн туршилтын орчин болгон ашиглаж уламжлалт JSF болон орчин үеийн React (бусад JavaScript фреймворкууд)  технологийг хослуулсан холимог микрофронтэнд архитектурын туршилтын загвар боловсруулж, түүнийгээ үнэлэх зорилго тавин ажиллаж байна. 

\hspace*{1cm}\textbf{Одоогоор хийж гүйцэтгэсэн ажлууд:}
\begin{itemize}
    \item Микрофронтэнд архитектур болон түүнийг хэрэгжүүлэх аргуудын талаарх онол, холбогдох судалгаануудыг судалж, туршилтын орчны бизнес процессуудыг амжилттай загварчиллаа.
    \item Системийн хэрэглэгчийн интерфейсийн зохиомжийг гаргаж, түүний дагуу React фрэймворк дээр Оюутан болон Багшийн модулиудын суурь хөгжүүлэлтийг хийж гүйцэтгэв.
    \item Хөгжүүлсэн React модулиудыг JSF хост аппликэйшн дотор Webpack хэрэгсэл болон скрипт тарилт ашиглан нэгтгэж, холимог микрофронтэнд архитектур бүхий туршилтын загварыг бодитоор хэрэгжүүлээд байна.
\end{itemize}
\hspace*{1cm}\textbf{Цаашид хийж гүйцэтгэх ажлууд:}
\begin{itemize}
    \item Хэрэгжүүлсэн туршилтын загвар дээр тулгуурлан ISO/IEC 25010 стандартын дагуу системийн чанарын үзүүлэлтүүдийг тодорхойлж, туршлагатай хөгжүүлэгчид болон багш нараас үнэлгээний санал асуулга цуглуулах;
    \item Цуглуулсан өгөгдлөөр DEMATEL математик аргачлалыг ашиглан тооцоолол хийж, чанарын шалгуур үзүүлэлтүүдийн хоорондын шалтгаан-үр дагаварын хамаарлыг тогтоох;
    \item Гарсан үр дүнд үндэслэн JSF болон React хосолсон архитектурын давуу болон сул талыг шинжлэх ухааны үндэслэлтэйгээр дүгнэж, дипломын эцсийн хамгаалалтад бэлтгэнэ.
\end{itemize}